\chapter{KIỂM THỬ, DEBUG VÀ QUAN SÁT}
\label{chap:testing}

Một chatbot hoạt động theo kiểu ``thỉnh thoảng trả lời có vẻ nhớ'' hoàn toàn không đủ tiêu chuẩn để đưa lên môi trường chạy thực tế. Chúng ta cần thiết lập các bài kiểm thử tự động để bảo đảm các luật bất biến (invariants) về bộ nhớ được duy trì xuyên suốt các đợt cập nhật tính năng. Nguyên tắc cốt lõi là việc kiểm thử bộ nhớ phải được tách biệt hoàn toàn khỏi sự bất định (non-determinism) của LLM bằng cách sử dụng các mô hình giả lập hoạt động theo luật cố định (rule-based mock models).

\section{Bốn quy tắc bất biến cốt lõi về bộ nhớ}

Mọi hệ thống bộ nhớ chatbot đều phải vượt qua bốn quy tắc kiểm thử sau:
\begin{enumerate}
  \item \textbf{Cùng luồng (Same Thread):} Tin nhắn sau bắt buộc phải đọc được thông tin ngữ cảnh của tin nhắn trước.
  \item \textbf{Khác luồng (Different Thread):} Trạng thái ngắn hạn (short-term memory) tuyệt đối không được phép rò rỉ qua lại giữa hai luồng trò chuyện khác nhau.
  \item \textbf{Cùng người dùng, khác luồng (Same User, Cross Thread):} Bộ nhớ dài hạn (long-term memory) trong Store phải có khả năng chia sẻ xuyên suốt các luồng của cùng một tài khoản.
  \item \textbf{Khác người dùng (Different User):} Bộ nhớ dài hạn của người dùng A tuyệt đối không được phép rò rỉ sang người dùng B.
\end{enumerate}

Dưới đây là mã nguồn kiểm thử tự động bốn quy tắc bất biến trên sử dụng thư viện \texttt{pytest}:

\begingroup
\lstinputlisting[inputencoding=utf8/latin1,language=Python,caption={Test short-term và long-term isolation}]{code/tests/test_memory.py}
\endgroup

\subsection{Phân tích chi tiết mã nguồn \texttt{test\_memory.py}}
\begin{itemize}
  \item \textbf{Hàm \texttt{test\_two\_threads\_do\_not\_mix\_names}:} Khởi tạo đồ thị bằng mô hình luật xác định, tạo hai luồng ngẫu nhiên bằng UUID. Người dùng A nói tên Nhật ở luồng A, người dùng B nói tên An ở luồng B. Hàm assert kiểm tra kết quả trả về của từng luồng độc lập có chính xác hay không.
  \item \textbf{Hàm \texttt{test\_profile\_is\_separated\_by\_user\_id}:} Lưu trữ thông tin trực tiếp vào Store dưới hai tài khoản khác nhau (\texttt{user-a} và \texttt{user-b}) và thực hiện assert kiểm tra sự cách ly dữ liệu ở tầng Store vật lý.
  \item \textbf{Hàm \texttt{test\_long\_term\_memory\_crosses\_threads\_for\_same\_user}:} Khởi tạo dịch vụ \texttt{ChatService}. Gửi câu lệnh ghi nhớ ở một luồng chat ngẫu nhiên của người dùng Nhật, sau đó thực hiện câu hỏi hỏi tên ở một luồng chat mới tinh của cùng người dùng đó. Kết quả assert chứng minh bộ nhớ dài hạn đã hoạt động xuyên luồng hội thoại.
  \item \textbf{Hàm \texttt{test\_long\_term\_memory\_does\_not\_cross\_users}:} Tương tự như trên nhưng câu hỏi được thực hiện bởi tài khoản \texttt{user-b}. Kết quả assert chứng minh tài khoản B nhận được câu trả lời mặc định là chưa biết tên, không bị ảnh hưởng bởi dữ liệu của tài khoản A.
\end{itemize}

\section{Kiểm thử tầng API, bảo vệ ranh giới bảo mật}

Đồ thị LangGraph chạy độc lập có thể đúng về mặt logic bộ nhớ, nhưng khi đóng gói vào API của FastAPI, hệ thống vẫn có khả năng bị rò rỉ thông tin nếu các lập trình viên quên thực hiện việc kiểm tra quyền sở hữu tại các hàm route. Do đó, viết kiểm thử tích hợp (integration tests) cho API bằng \texttt{TestClient} là cực kỳ quan trọng:

\begingroup
\lstinputlisting[inputencoding=utf8/latin1,language=Python,caption={Test ownership tại ranh giới FastAPI}]{code/tests/test_fastapi_app.py}
\endgroup

\subsection{Phân tích chi tiết mã nguồn \texttt{test\_fastapi\_app.py}}
\begin{itemize}
  \item \textbf{Hàm \texttt{test\_conversation\_flow\_and\_ownership}:} Giả lập một client gọi endpoint POST tạo hội thoại cho \texttt{user-a}. Sau đó, giả lập một client đại diện cho tài khoản khác (\texttt{user-b}) cố tình gọi endpoint GET để đọc danh sách tin nhắn của hội thoại vừa tạo. Kết quả assert kiểm tra mã trạng thái trả về có đúng là lỗi HTTP 403 (Forbidden) hay không.
  \item \textbf{Hàm \texttt{test\_unknown\_thread\_returns\_404}:} Kiểm tra hành vi hệ thống khi người dùng gọi một \texttt{thread\_id} không tồn tại trên hệ thống. Kết quả assert xác nhận lỗi HTTP 404 (Not Found) được trả về chính xác.
\end{itemize}

\section{Failure Lab: Các lỗi kinh điển và cách xử lý}

Hãy chủ động làm quen với các tình huống lỗi sau trong lúc xây dựng hệ thống:
\begin{itemize}
  \item \textbf{Sử dụng thread ID tĩnh (Static Thread ID):} Gán cứng một giá trị \texttt{thread\_id} trong code client. Triệu chứng: Mọi người dùng trên hệ thống đều nhìn thấy tin nhắn của nhau và phản hồi của trợ lý bị lẫn lộn hoàn toàn.
  \item \textbf{Gửi trùng lặp lịch sử từ Client (Duplicate History):} Client vừa gửi lịch sử hội thoại vừa sử dụng checkpointer. Triệu chứng: Số lượng token đầu vào tăng vọt nhanh chóng, câu trả lời của mô hình lặp đi lặp lại do reducer nhân đôi tin nhắn cũ.
  \item \textbf{Lưu trữ tạm thời bị mất dữ liệu (Volatile Storage):} Sử dụng saver dạng RAM tạm thời trong môi trường Docker container. Triệu chứng: Dữ liệu hội thoại biến mất hoàn toàn mỗi khi container bị crash hoặc cập nhật phiên bản mới.
  \item \textbf{Namespace Store thiếu user\_id (Missing User Namespace):} Lưu trữ hồ sơ người dùng vào Store mà không phân vùng namespace theo \texttt{user\_id}. Triệu chứng: Người dùng đăng nhập sau ghi đè tên của người dùng đăng nhập trước trên toàn hệ thống.
\end{itemize}

\section{Giám sát Vận hành (Telemetry \& Metrics)}

Khi đưa hệ thống bộ nhớ vào vận hành thực tế, chúng ta cần theo dõi sát sao các chỉ số (metrics) sau để đảm bảo hệ thống chạy mượt mà và tối ưu hóa chi phí:

\begin{table}[H]
\centering
\caption{Các chỉ số giám sát vận hành bộ nhớ tác tử}
\begin{tabular}{p{5.8cm}p{4.0cm}p{4.8cm}}
\toprule
\textbf{Metric cần đo đạc} & \textbf{Mục đích giám sát} & \textbf{Cảnh báo cần lưu ý} \\
\midrule
Mã băm (hash) của thread và user ID & Tra vết luồng yêu cầu (traceability) & Tránh log các thông tin thô dạng PII của người dùng. \\
Số lượng tin nhắn trong checkpoint & Phát hiện lỗi nhân đôi lịch sử trò chuyện & Số lượng tin nhắn tăng bất thường sau mỗi lượt chat. \\
Tổng số token đầu vào/đầu ra (mỗi provider) & Kiểm soát chi phí tài chính & Token đầu vào vượt quá Token budget đã thiết lập. \\
Độ trễ đọc/ghi checkpoint và Store & Đánh giá hiệu năng database & Độ trễ tăng cao do nghẽn kết nối DB hoặc thiếu index. \\
Tần suất chạy tóm tắt (summary trigger) & Tối ưu hóa hiệu năng gọi LLM & Tóm tắt chạy liên tục sau mỗi lượt gọi (tốn token). \\
Số lượng mã lỗi bảo mật HTTP 403 & Cảnh báo an ninh/lỗi giao diện client & Tần suất lỗi 403 đột biến báo hiệu hành vi tấn công dò quét. \\
\bottomrule
\end{tabular}
\end{table}

\begin{aicheckpoint}[Lab 20 phút]
Hãy chạy toàn bộ các bài test tự động bằng lệnh \texttt{uv run pytest}. Sau đó, thử cố tình loại bỏ tham số \texttt{user\_id} khỏi tuple namespace trong file \texttt{long\_term\_memory.py} rồi chạy lại pytest. Hãy quan sát xem bài test nào bị thất bại và giải thích cơ chế phát hiện lỗi bảo mật của bài test đó.
\end{aicheckpoint}
